\section{Bethlehem, the Manger, and the Holy Night}\label{sec:bethlehem}

With First Vespers of the Nativity on the evening of 24 December, Advent yields
to Christmas. All three Masses of the Nativity retain their ancient Roman
stations in the 1962 Missal: the First Mass at Midnight at St.~Mary Major
\textit{ad Praesepe}---``at the Crib''; the Second Mass at Dawn at
St.~Anastasia; and the Third Mass in the Day at St.~Mary Major. Here, on the
threshold of the holy night, every thread that the Church has held through four
weeks of watching comes home to a feeding trough in a small town of Judea.

\subsection{Bethlehem, the House of Bread}

On Christmas Day at St.~Mary Major, Pope St.~Gregory the Great apologized to
his people for preaching briefly---``because, by the Lord's bounty, we are
about to celebrate the solemnities of Mass three times today''---and then went
straight to the name of the town where Mary laid her child. In Homily 8 on the
Gospels, read at Christmas Matins, he asks why Caesar Augustus enrolled the
whole world at that hour (to show that One was coming in the flesh who would
enroll his elect in eternity) and why the Child was born in Bethlehem:

\begin{quote}
Then, the Lord is born in Bethlehem. Now the name Bethlehem signifieth ``the
House of Bread,'' and thus it is the birth-place of Him Who hath said, ``I am
the Living Bread, Which came down from heaven.'' (John vi. 51.) We see then
that this name of Bethlehem was prophetically given to the place where Christ
was born, because it was there that He was to appear in the flesh by Whom the
souls of the faithful are fed unto life eternal. He was born, not in His
Mother's house, but away from home. And this is a mystery, showing that this
our mortality into which He was born was not the home of Him Who is begotten of
the Father before the worlds.
\end{quote}

In the next sentences of the same homily, Gregory looks straight into the
manger. Isaiah had cried in Advent, ``All flesh is grass'' (Is 40:6; Latin
\textit{fenum}, hay), and Jesus will say before his Passion, ``Unless the grain
of wheat falling into the ground die, itself remaineth alone'' (John 12:24).
``Made man,'' Gregory says, ``he turned our hay into wheat'' (\textit{fenum
nostrum vertit in frumentum}), and therefore ``when born he is laid in a
manger, that he might refresh all the faithful, as sacred animals, with the
wheat of his flesh, lest they remain fasting from the food of eternal
understanding.''

East and West heard the same wonder in that name. In Antioch, St.~John
Chrysostom calls his people to ``hasten to Bethlehem, where is the house of the
spiritual Bread,'' naming the Eucharistic altar ``the spiritual manger''
(Homily~7 on Matthew). In Alexandria, St.~Cyril (\textit{Catena Aurea} on Luke
2:7) sees Christ laid in the feeding trough to lift men from beastly appetite
to heavenly food, ``partaking not of hay, but of the heavenly bread, the
lifegiving body.''

St.~Bede turns that geography into the daily life of every believer: ``down to
the very end of time, the Lord ceases not to be conceived at Nazareth, to be
born at Bethlehem, whenever any of His hearers taking of the flour of the word
makes himself a house of eternal bread.'' Preaching on Christmas Eve,
St.~Bernard of Clairvaux carries the image straight to Holy Communion:
``Bethlehem is the house of bread; Juda signifies confession or praise. If,
then, you replenish your soul with the food of the Divine Word, the Body of our
Lord Jesus Christ, and devoutly receive the Bread which came down from heaven,
and which giveth life to the world... then, indeed, you are become a Bethlehem
fitted to receive our Lord.''

Now the verse that opened the first night of Advent---Isaiah 1:3, ``The ox
knoweth his owner, and the ass his master's crib''---arrives at its
fulfilment. As St.~Gregory Nazianzen cries in Oration~38.17, ``honour little
Bethlehem, which hath led thee back to Paradise; and worship the manger through
which thou, being without sense, wast fed by the Word.'' At Christmas Matins,
after the fourth lesson---the first from St.~Leo---the Church sings her Fourth
Responsory:

\begin{massbox}{Fourth Responsory of Christmas Matins: \textit{O magnum mysterium}}
\textbf{Latin.} \textbf{R.} O magnum mysterium, et admirabile sacramentum, ut
animalia viderent Dominum natum, jacentem in praesepio: * Beata Virgo, cujus
viscera meruerunt portare Dominum Christum. \textbf{V.} Ave Maria, gratia
plena: Dominus tecum. Beata.

\textbf{English.} How great is this mystery, how wonderful is the teaching of
the faith! The beasts saw the new-born Lord lying in a manger. Blessed is that
Virgin whose womb was made meet to bear our Lord Christ. \textbf{Verse.} Hail,
Mary, full of grace: the Lord is with thee. \textbf{Answer.} Blessed is that
Virgin whose womb was made meet to bear our Lord Christ.
\end{massbox}

\subsection{Wrapped and Laid: The Manger and the Tomb}

St.~Luke tells the birth of Jesus in a single unhurried sentence (Luke 2:7):

\begin{quote}
And she brought forth her first born son and wrapped him up in swaddling
clothes and laid him in a manger: because there was no room for them in the
inn.
\end{quote}

At the close of Good Friday, St.~Luke places Joseph of Arimathea beside the
cross with the same unhurried plainness (Luke 23:53):

\begin{quote}
And taking him down, he wrapped him in fine linen and laid him in a sepulchre
that was hewed in stone, wherein never yet any man had been laid.
\end{quote}

In St.~Luke's Greek, each scene keeps its own proper verbs: at Bethlehem Mary
swathes the child in swaddling bands (\textit{esparganosen}) and reclines him
(\textit{aneklinen}) in a manger; on Good Friday Joseph of Arimathea winds the
body (\textit{enetylixen}) in linen and places him (\textit{etheken}) in a
rock-hewn tomb. The Latin Vulgate renders the first gesture with
\textit{involvit} in both verses while keeping \textit{reclinavit} and
\textit{posuit} distinct, and the Douay-Rheims translates both actions with
``wrapped'' and ``laid.'' The two scenes mirror each other in deed: a faithful
hand wrapping the body of the Lord in linen and laying him in a hollowed place.

The Fathers held those two cloths together in a single gaze. St.~Gregory
Nazianzen sets the birth and the Resurrection side by side in his Third
Theological Oration (Oration 29.19):

\begin{quote}
He dwelt in the womb---but He was recognized by the Prophet, himself still in
the womb, leaping before the Word, for Whose sake He came into being. He was
wrapped in swaddling clothes---but He took off the swathing bands of the grave
by His rising again. He was laid in a manger---but He was glorified by Angels,
and proclaimed by a star, and worshipped by the Magi.
\end{quote}

At Bethlehem, St.~Jerome observes in \textit{Contra Helvidium} (preserved in
the \textit{Catena Aurea} on Luke 2:7), ``here was no midwife, no tender
anxiety of women; she wrapped the Child up in swaddling clothes, herself both
mother and midwife.'' St.~Bede reads the binding of those tiny hands and feet
as our release: ``He who clothes the whole world with its varied beauty, is
wrapped up in common linen, that we might be able to receive the best robe; He
by Whom all things are made, is folded both hands and feet, that our hands
might be raised up for every good work, and our feet directed in the way of
peace.'' Even the Magi's gifts carry the shadow of Friday inside the joy of
Bethlehem: ``gold, frankincense, and myrrh'' (Matt 2:11)---myrrh, as
St.~Gregory Nazianzen noted (Oration 38.17), offered to One who dies for us.
St.~Thomas Aquinas notes in the \textit{Summa Theologiae} (III, q.~35, a.~7)
why the Lord chose Bethlehem to be born and Jerusalem to suffer: ``Christ, on
the contrary, willed to be born in a mean city, and to suffer reproach in a
great city.'' And St.~Maximus of Turin tells any heart tempted to look away
from such poverty: ``If thou despisest the manger, raise thy eyes a little, and
behold the new star in heaven proclaiming to the world the Lord's nativity.''

\subsection{The Wonderful Exchange}

Why did the eternal Son accept the bands of infancy, the feeding trough, and
the rock-hewn tomb? From the second century onward, Greek and Latin Fathers
answered with one startling sentence. St.~Irenaeus of Lyons writes in the
preface to Book V of \textit{Against Heresies} of ``our Lord Jesus Christ, who
did, through His transcendent love, become what we are, that He might bring us
to be even what He is Himself.'' St.~Athanasius of Alexandria puts it in ten
words in \textit{On the Incarnation} (54.3): ``For He was made man that we
might be made God.'' St.~Gregory Nazianzen marvels in Oration 38.13: ``And He
Who gives riches becomes poor, for He assumes the poverty of my flesh, that I
may assume the richness of His Godhead.''

In the Latin West, St.~Augustine wakes his congregation in Sermon 185.1 with a
thunderclap: \textit{Expergiscere, homo: pro te Deus factus est homo}---``Wake
up, O man: for you God was made man.'' You would have died forever had he not
been born in time; \textit{Perisses, nisi venisset}---``You would have perished,
had he not come.'' In Sermon 191.1, Augustine watches the Creator step inside
his own clock: \textit{Homo factus, hominis factor}---``Made man, the Maker of
man''---so that he who governs the stars might nurse at a mother's breast, the
Bread might hunger, the Fountain might thirst, and Life itself might die.
St.~Bernard catches the same wonder in his Second Homily on the \textit{Missus
est}: ``There we behold light withholding its rays, the Word an infant, the
Living Water athirst, Him Who is the Bread of Heaven suffering hunger.''

Pope St.~Leo the Great turns that exchange into a summons to every baptized
Christian in his First Sermon on the Nativity (Sermon 21.1, 3):

\begin{quote}
Our Saviour, dearly-beloved, was born to-day: let us be glad. For there is no
proper place for sadness, when we keep the birthday of the Life, which destroys
the fear of mortality and brings to us the joy of promised eternity...
Christian, acknowledge thy dignity, and becoming a partner in the Divine
nature, refuse to return to the old baseness by degenerate conduct. Remember
the Head and the Body of which thou art a member.
\end{quote}

And in Sermon 26.2, Leo binds Christ's birth straight to our baptism: ``For the
birth of Christ is the source of life for Christian folk, and the birthday of
the Head is the birthday of the body.'' The Church's own song for that bargain
of mercy stands at Lauds and Vespers on the Octave Day of Christmas, 1 January:

\begin{massbox}{First Antiphon of Lauds, Octave Day of Christmas: \textit{O admirabile commercium}}
\textbf{Latin.} O admirabile commercium: * Creator generis humani, animatum
corpus sumens, de Virgine nasci dignatus est: et procedens homo sine semine,
largitus est nobis suam Deitatem.

\textbf{English.} How wonderful is the dealing here! * The Maker of man hath
taken to Himself a living body, and hath been pleased to be born of the Virgin:
He is come forth conceived without seed, and hath made us partakers of His
Divine nature.
\end{massbox}

\subsection{The Holy Night}

The grace of Bethlehem did not end in the first century. On Christmas Eve 1886,
after Midnight Mass in Lisieux, the young Th\'er\`ese Martin---still fragile
and easily reduced to tears since her mother's death---overheard a tired remark
from her father and felt an instantaneous, lasting strength enter her heart. In
her memoir she called it her complete conversion: ``On this night of grace, the
third period of my life began---the most beautiful of all, the one most filled
with heavenly favours. In an instant Our Lord, satisfied with my good will,
accomplished the work I had not been able to do during all these years.'' More
than a century later, Pope Francis, in his apostolic letter \textit{Admirabile
signum} (2019, nos.~1 and 10), urges parents and grandparents to pass on the
habit of setting up the Nativity scene at home so that children and adults
alike may kneel and contemplate the love of God who became a child.

For four weeks the Church has held back the angels' song for this night
(Section~\ref{sec:history}). At Midnight Mass, after the Introit \textit{Dominus
dixit ad me} (``The Lord hath said to me: Thou art my son, this day have I
begotten thee,'' Ps~2:7) and the \textit{Kyrie}, the priest intones
\textit{Gloria in excelsis Deo}, opening the great hymn with the angels' words
heard by the shepherds in the dark fields of Bethlehem: ``Glory to God in the
highest: and on earth peace to men of good will'' (Luke 2:14). The Collect asks
that ``we may enjoy in heaven the delights of Him whose mystical light we have
known on earth.''

At dawn, the shepherds say one to another, ``Let us go over to Bethlehem and
let us see this word that is come to pass, which the Lord hath shewed to us''
(Luke 2:15). And St.~Gregory Nazianzen opens the morning in Oration~38.1 with
the only cry left to speak:

\begin{quote}
Christ is born, glorify ye Him. Christ from heaven, go ye out to meet Him.
Christ on earth; be ye exalted. Sing unto the Lord all the whole earth; and
that I may join both in one word, Let the heavens rejoice, and let the earth be
glad, for Him Who is of heaven and then of earth. Christ in the flesh, rejoice
with trembling and with joy; with trembling because of your sins, with joy
because of your hope.
\end{quote}
